\section{Certificates, obligations, clocks, and coupling}
\label{sec:certificates}

A restricted defender search is useful only when a verifier can reconstruct
why each omitted reply is harmless.  The certificate therefore records exact
game states, exact searched successors, the cells needed by future attacker
moves, and separate clocks for occupancy and completion hazards.  A second
play, called the ghost play, follows the recorded tree while the real defender
may choose an omitted reply.

\subsection{Exact certificate grammar}

The grammar fixes the data that a verifier receives.  In particular, a leaf
is not a bare tactical verdict.  It contains the windows and deadlines needed
to replay that verdict after real and ghost defender stones have diverged.

\begin{definition}[Certificate]
\label{def:D9}
A \emph{certificate} $\mathcal C$ for the assertion that $\A$ wins from
$P_0$ is a finite rooted tree.  Every node is labelled by an exact position
in the sense of \cref{def:D3}, including its mover and remaining placement
budget.  The root is the nonterminal position $P_0$.  The verifier recomputes
the ply clock as the root ply plus path depth.  The \emph{leaf-ply} is the
index of the last placement on the path into a leaf.  Every edge is the exact
legal successor specified by \cref{def:D4}.  No edge is defender-terminal.

An \emph{OR-COMPLETION leaf} has $\A$ to place.  It names one designated
placement, an $\A$-complete witness window for that placement, and the
completion ply.

A \emph{WIN leaf} has $\A$ to place and passes a verifier check of
$\ownwinnow$.  It names an $\A$-alive count-$5$ witness for either budget, or
an $\A$-alive count-$4$ witness when the attacker budget is $2$.  Its declared
resolution lies in the current attacker turn.

A \emph{LOSS leaf} has $\D$ to place with budget $b$.  It names a family
$\mathcal T$ of $\A$-threat windows whose empty-set family has transversal
number $\tau(\mathcal T)>b$, and it passes the defender check
$\neg\ownwinnow$.  For every complete nonterminal remainder $H$ of the
defender turn, consisting of exactly $b$ placements unless the game ends
earlier, some $W_H\in\mathcal T$ satisfies
\[
  E(W_H,P_L)\cap H=\varnothing.
\]
The attacker then fills the at most two remaining empties of $W_H$ during the
following turn.  The declared worst-case resolution is
\[
  \text{leaf-ply}+b+2.
\]
The verifier may require at most three named windows when $b=1$ and at most
five when $b=2$, by \cref{lem:L13}.

An internal OR node has one designated attacker placement and its exact
child.  An internal AND node $N$ has an independently nonempty set $S(N)$ of
searched legal defender replies and one exact child for each member of
$S(N)$.  Every other legal defender reply carries a dismissal claim.  The
verifier retains the internal defender check $\neg\ownwinnow$ as a diagnostic.
It follows from the ordinary completion zone by \cref{lem:L14}; it remains
mandatory at LOSS leaves and is an explicit premise of \cref{thm:T6}.  Every
terminal resolution is finite.  The global horizon $T$ is the maximum of the
declared resolutions on all root-to-leaf paths.  The finite DAG form is given
in \cref{def:D18}.
\end{definition}

\begin{plainterms}
The certificate is an exact, finite claim about one attacker strategy.  OR
nodes state what the attacker will do.  AND nodes retain enough defender
replies to justify every omitted one.  The typed leaves prevent a verifier
from accepting a tactical label without the cells and time needed to realize
it.
\end{plainterms}

\subsection{A certificate used throughout the paper}
\label{subsec:running-example}

One finite certificate supplies the numerical example in this and the next
three sections.  Its root has $\D$ to place with budget $1$.  There are
fifteen stones of each colour:
\begin{align*}
A_0={}&\{(q,0):0\le q\le4\}\cup\{(0,-7)\}\\
 &\cup\{(-3,20),(-2,20),(-1,20),
          (0,17),(0,18),(0,19),\\[-1mm]
 &\hspace{32mm}(-3,23),(-2,22),(-1,21)\},\\
D_0={}&\{(-1,0),(8,0),(8,8),(8,16),(8,24),(8,32),\\[-1mm]
 &\hspace{9mm}(0,40),(1,40),(2,40),
          (-8,0),(0,-8),(-8,8),\\[-1mm]
 &\hspace{32mm}(8,-8),(16,-8),(16,0)\}.
\end{align*}
The support chain makes every displayed module legal under the radius-$8$
rule.  Enumeration of the 18 incident windows per stone finds no complete
window and finds at most three defender stones in any defender-alive window
\cite{hexobot-companions}.

The root has the singleton attacker threat
\[
  T=\{(0,0),(1,0),(2,0),(3,0),(4,0),(5,0)\},
  \qquad h=(5,0).
\]
Its empty-set family is $\{\{h\}\}$.  The exact searched reply is $h$.
The attacker then places at $(3,20)$ and $(0,20)$ and reaches a defender
budget-$2$ LOSS leaf.  The three witness empty sets there are
\[
 \{(1,20),(2,20)\},\qquad
 \{(0,21),(0,22)\},\qquad
 \{(1,19),(2,18)\}.
\]
They are pairwise disjoint, so their transversal number is $3$.

\begin{figure}[tbp]
  \centering
  \begin{subfigure}[t]{.31\textwidth}
    \centering
    \begin{tikzpicture}[scale=.235]
      \HexPatch{-1/-1,0/-1,1/-1,2/-1,3/-1,4/-1,5/-1,6/-1,7/-1,8/-1,9/-1,10/-1,
        -1/0,0/0,1/0,2/0,3/0,4/0,5/0,6/0,7/0,8/0,9/0,10/0,
        -1/1,0/1,1/1,2/1,3/1,4/1,5/1,6/1,7/1,8/1,9/1,10/1}
      \HexWindow[window one]{0}{0}{1}{0}
      \HexWindow[window two]{5}{0}{1}{0}
      \foreach \q in {0,...,4}{\HexAttacker{\q}{0}}
      \HexDefender{-1}{0}\HexDefender{8}{0}
      \HexMark{5}{0}{cross}{}
      \node[font=\scriptsize,above=2pt] at \HexAt{5}{0} {$h$};
    \end{tikzpicture}
    \caption{The singleton threat $T$.  The outlined cells form the later
      queried window $W'$.}
    \label{fig:running-gate-module}
  \end{subfigure}\hfill
  \begin{subfigure}[t]{.35\textwidth}
    \centering
    \begin{tikzpicture}[scale=.27]
      \HexRhombus{-3}{3}{-3}{3}
      \foreach \q/\r in {3/0,0/0,1/0,2/0,0/1,0/2,1/-1,2/-2}{
        \HexShadeCell[zone direct]{\q}{\r}}
      \foreach \q/\r in {-3/0,-2/0,-1/0,0/-3,0/-2,0/-1,-3/3,-2/2,-1/1}{
        \HexAttacker{\q}{\r}}
      \HexMark{3}{0}{dot}{}\HexMark{0}{0}{dot}{}
      \node[font=\scriptsize,above] at \HexAt{3}{0} {$(3,20)$};
      \node[font=\scriptsize,below right] at \HexAt{0}{0} {$(0,20)$};
    \end{tikzpicture}
    \caption{The LOSS module, translated by $(0,-20)$ for display.  The
      shaded cells are the eight future cells named by this certificate
      fragment; the direct-zone style previews their role at an ordinary
      zone node.}
    \label{fig:running-loss-module}
  \end{subfigure}\hfill
  \begin{subfigure}[t]{.28\textwidth}
    \centering
    \begin{tikzpicture}[scale=.245]
      \HexRhombus{-4}{6}{-1}{1}
      \HexWindow[window one]{0}{0}{1}{0}
      \foreach \q in {-3,-2,-1,3,4,5}{\HexShadeCell[zone touched]{\q}{0}}
      \foreach \q in {0,1,2}{\HexDefender{\q}{0}}
      \node[font=\scriptsize,above] at \HexAt{1}{1}
        {$V_{-3},V_{-2},V_{-1},V_0$};
    \end{tikzpicture}
    \caption{The four qualifying windows at height $40$, translated by
      $(0,-40)$.  Their six shaded empties are the complete ordinary touched
      term at exposure $3$.}
    \label{fig:running-window-module}
  \end{subfigure}

  \vspace{1ex}
  \begin{tabular}{@{}cl@{\qquad}cl@{}}
    \HexLegendSwatch[zone direct] & direct-role annotation &
    \HexLegendSwatch[zone seed] & ranked seed guard: $\varnothing$\\
    \HexLegendSwatch[zone touched] & touched-window annotation &
    \HexLegendSwatch[zone virgin] & virgin-window guard: $\varnothing$
  \end{tabular}
  \caption{Tactical modules of the 15--15 running position.  Blank axial
  distance between the modules is suppressed; the coordinate lists in the
  text specify the complete position.  The direct and touched shadings give
  the exact ordinary full-clock terms; the key records that the seed and
  virgin terms are empty.  The root is a forcing gate and instead searches
  the exact kernel $\{h\}$; the complete ordinary decomposition and the
  verified exposure calculations for $W$ and $W'$ appear in
  \cref{ex:running-zone-computation}.  Filled disks are attacker stones and
  open disks are defender stones.}
  \label{fig:running-position}
\end{figure}

\begin{figure}[tbp]
  \centering
  \begin{tikzpicture}[
      cert/.style={draw,rounded corners=2pt,align=center,minimum height=8mm,
        inner sep=4pt,font=\small},
      every edge/.style={draw,-{Stealth[length=5pt]}}]
    \node[cert] (q) {$P_0$\\$\D,b=1$};
    \node[cert,right=18mm of q] (a2) {$\A,b=2$};
    \node[cert,right=18mm of a2] (a1) {$\A,b=1$};
    \node[cert,right=18mm of a1] (loss) {LOSS\\$\D,b=2$, $\tau=3$};
    \node[cert,below=13mm of loss] (win) {$\A$ completes\\a surviving witness};
    \path (q) edge node[above,font=\scriptsize] {$h=(5,0)$} (a2)
          (a2) edge node[above,font=\scriptsize] {$(3,20)$} (a1)
          (a1) edge node[above,font=\scriptsize] {$(0,20)$} (loss)
          (loss) edge node[right,font=\scriptsize] {$b+2=4$ plies} (win);
  \end{tikzpicture}
  \caption{The certificate fragment used in the running example.  The root
  has one exact searched reply.  The two attacker OR edges lead to a LOSS
  leaf with three disjoint two-cell witnesses.  Section~\ref{sec:gates}
  gives the forcing-gate interpretation of the root edge.}
  \label{fig:running-certificate}
\end{figure}

\subsection{Obligations and local clocks}

A global ply horizon overstates the work remaining on most certificate
branches.  Three local objects replace it: a defender-placement budget for a
whole subtree, a deadline for each needed cell, and an exposure for each
window.

\begin{definition}[Admissible local defender budget]
\label{def:D14}
The exact local budget $B(N)$ is measured in defender placements and is
defined by
\[
\begin{aligned}
B(\text{OR-COMPLETION})&=B(\text{WIN})=0,
&B(\text{LOSS},b)&=b,\\
B(\text{OR})&=B(\text{child}),
&B(\text{AND})&=1+\max_C B(C).
\end{aligned}
\]
The two placements of a defender turn are two D4 edges.  A LOSS leaf counts
its remaining $b$ placements.  Exact maxima are optional: a verifier may
accept any nonnegative integral labels satisfying $B(L)\ge b$ at every LOSS
leaf, $B(N)\ge B(C)$ on every OR edge, and
$B(N)\ge1+B(C)$ on every AND edge.
\end{definition}

\begin{plainterms}
The budget counts only defender placements that can still occur on the
selected branch.  An attacker edge costs nothing.  A defender edge costs one,
and a LOSS leaf reserves enough budget for the remainder that is not expanded
as certificate edges.
\end{plainterms}

Only cells that the attacker may later need must be protected from a
dismissed defender stone.  The obligation definition records those uses and
their exact deadlines.

\begin{definition}[Compressed obligations]
\label{def:D10}
Write $E(W,P)=W\setminus\Stones(P)$.  A \emph{live obligation role} below a
node $N$ is either a future designated certificate attacker placement,
including an OR-COMPLETION placement, or a triple $(L,W,y)$ for a named WIN
or LOSS witness $W$ at a leaf $L$ and a cell $y\in E(W,P_L)$.  The second kind
represents every permitted leaf continuation placement, including the
adaptive LOSS choice.

Let $\Omega(N)$ be all such roles at reachable descendants of $N$.  Define
\[
  \Obl(\mathcal C,N)=\Prot(N)
  =\{y:\text{some live role in $\Omega(N)$ is carried by $y$}\}.
\]
An attacker-placement role remains live through the check immediately before
its designated OR step.  A witness-empty role remains live through the leaf
entry check.  Each role is discharged immediately after that check.  A full
witness window is not itself an obligation.  Reachable-descendant nesting
gives $\Prot(M)\subseteq\Prot(N)$ whenever $M$ descends from $N$.
\end{definition}

\begin{plainterms}
An obligation is a future use of one cell, not a permanent reservation of an
entire witness window.  A dismissed defender stone may alter unrelated cells
without invalidating the certificate.  The verifier protects each needed
cell only until the moment when that use is checked.
\end{plainterms}

A cell-specific clock charges only defender opportunities that occur before
one named role is checked.

\begin{definition}[Cell-specific deadlines]
\label{def:D15}
The deadline of a role $\rho$ is its designated attacker placement, or leaf
entry for a WIN or LOSS witness-empty role.  At a node $N$ from which $\rho$
is reachable, $r_N(\rho)$ is an integral upper bound on the number of defender
placements before that deadline.  Exact ranks count the maximum number of AND
edges on a reachable path.  While the role is live, the verifier checks
\[
 r_N(\rho)\ge r_C(\rho)\quad\text{on an OR edge},
 \qquad
 r_N(\rho)\ge1+r_C(\rho)\quad\text{on an AND edge},
\]
and checks $r=0$ at the deadline.  If several roles use one cell, set
\[
  r_N(y)=\max\{r_N(\rho):\rho\in\Omega(N)
                       \text{ is live and carried by }y\}.
\]
The role is discharged only after its deadline check.  A seed band is formed
at an internal AND node only when $r_N(y)\ge1$; no negative radius is formed
at $r=0$.  LOSS witness empties need protection only through leaf entry,
because the adaptive leaf contract controls the remaining placements.
Defender-completion windows use \cref{def:D14,def:D16}, not role ranks.  The
uniform admissible assignment before a deadline is $r_N(\rho)=B(N)$.
\end{definition}

A separate window clock stops when the planned attacker first enters that
window, because later defender fills cannot precede that entry.

\begin{definition}[Per-window defender exposure]
\label{def:D16}
For each window $W$, let $E_Q^\D(W)$ be the maximum number of future defender
placements before the certificate attacker wins or first places in $W$.  Its
exact recurrence is
\[
E_Q^\D(W)=
\begin{cases}
0,&Q\text{ is WIN or OR-COMPLETION},\\
b,&Q\text{ is LOSS with defender budget }b,\\
0,&Q\text{ is an OR whose designated move enters }W,\\
E_C^\D(W),&Q\text{ is any other ordinary OR},\\
1+\max_C E_C^\D(W),&Q\text{ is an internal AND}.
\end{cases}
\]
If $W$ is already not alive for $\D$, set $E_Q^\D(W)=0$.  The attacker
\emph{enters} $W$ when its placement makes $W$ permanently not alive for
$\D$; the window need not become dead.  Each overlapping window has its own
clock.  The designated stopping move in $W$ is also an obligation, so a
separate occupancy check keeps that stop playable in the real game.
\end{definition}

\begin{plainterms}
Role ranks measure how many opportunities remain to occupy a particular
needed cell.  Window exposure measures how many defender stones can still
enter one particular window before the attacker blocks it or wins.  These are
different hazards, so one scalar is not used for both.
\end{plainterms}

The first few threat-creating attacker moves can be located without examining
their eventual certificate nodes.

\begin{lemma}[Short-range attacker-obligation containment]
\label{lem:L10}
In a certificate whose attacker placements are all threat-creating, let a
placement be the attacker's $k$th future placement from $N$.  If $k\le3$,
then its threat window contains at least
\[
  3-(k-1)\ge1
\]
attacker stones already present in $P_N$.  Its obligation therefore lies in
an attacker-touched window.  From the fourth future placement onward, a setup
cell may lie in a window that is virgin at $N$.
\end{lemma}

\begin{proof}[Proof sketch]
Immediately before a threat-creating placement, its witness window has at
least three attacker stones.  At most $k-1$ of them can be earlier future
placements, leaving the displayed number in $P_N$.  The fourth placement is
the first for which this count may be zero.  See
\cite{hexobot-companions}.
\end{proof}

The recurrences provide the inequalities used when the coupling descends to a
certificate child.

\begin{lemma}[Local-clock and nesting facts]
\label{lem:L11}
If a descendant $M$ is reached from $N$ after $k$ defender placements, then
\[
\begin{gathered}
B(M)+k\le B(N),\qquad
\cnt_\D(W,P_M)+B(M)\le\cnt_\D(W,P_N)+B(N),\\
r_M(\rho)+k\le r_N(\rho)
\quad\text{for every still-live role }\rho.
\end{gathered}
\]
Before the attacker enters $W$,
\[
 E_M^\D(W)+k\le E_N^\D(W),\qquad
 \cnt_\D(W,P_M)+E_M^\D(W)
 \le\cnt_\D(W,P_N)+E_N^\D(W).
\]
Every selected path below $N$, including a LOSS remainder, has at most
$B(N)$ defender placements before its declared attacker resolution.  Also
$E_N^\D(W)\le B(N)$, every exact role rank is at most $B(N)$, and
$\Prot(M)\subseteq\Prot(N)$.
\end{lemma}

\begin{proof}[Proof sketch]
Sum the verifier inequalities along the path.  Each AND edge contributes one
and each OR edge zero; the LOSS base covers its unexpanded remainder.  A
window gains at most one defender stone on each counted defender edge, and
reachable roles can only disappear after descent.  The D16 recurrence gives
the corresponding fixed-window bounds until its stated stop.  See
\cite{hexobot-companions}.
\end{proof}

\subsection{The ranked zone}

The searched set must block three distinct failure modes.  A defender may
occupy a cell needed later, may approach such a cell through a chain of newly
legal moves, or may complete a window before the mapped attacker blocks it.
Touched and virgin windows require different completion tests.

For a set $U$ of cells, write
\[
 B_s(U)=\{c:\dist(c,u)\le s\text{ for some }u\in U\},
 \qquad
 \dist(c,W)=\min_{w\in W}\dist(c,w).
\]

\begin{definition}[Zone-carrying certificate]
\label{def:D11}
At every internal AND node $N$, define
\begin{align*}
\Zdir(N)={}&\Prot(N)\cap\Legal(P_N),\\
\Zseed(N)={}&
 \bigcup_{\substack{y\in\Prot(N)\setminus
   (\Legal(P_N)\cup\Stones(P_N))\\ r_N(y)\ge1}}
 \bigl(\Legal(P_N)\cap B_{8(r_N(y)-1)}(\{y\})\bigr),\\
\Ztouch(N)={}&
 \bigcup\left\{E(W,P_N):
 \begin{array}{l}
 W\text{ is alive for }\D,\ \cnt_\D(W,P_N)\ge1,\\[-1mm]
 \cnt_\D(W,P_N)+E_N^\D(W)\ge6
 \end{array}\right\},\\
\Zvirgin(N)={}&
 \left\{c\in\Legal(P_N):
 \begin{array}{l}
 \text{some all-empty window }W\text{ has }E_N^\D(W)\ge6,\\[-1mm]
 \dist(c,W)\le8(E_N^\D(W)-6)
 \end{array}\right\}.
\end{align*}
Every cell in $\Ztouch(N)$ is legal: \cref{lem:L1} places it within distance
at most $5$ of an in-window defender stone.

The certificate is \emph{zone-carrying} when $S(N)$ is independently
nonempty at every internal AND node and the following verifier clauses hold.
\begin{enumerate}[label=\textup{(Z\arabic*)},leftmargin=12mm]
  \setcounter{enumi}{1}
  \item \label{zone:Z2}
  $S(N)\supseteq\Zdir(N)\cup\Ztouch(N)\cup\Zvirgin(N)$.
  \setcounter{enumi}{3}
  \item \label{zone:Z4}
  Every certificate attacker placement, including one permitted by a leaf
  contract, is within distance $8$ of an attacker stone in its predecessor
  position or a root stone unaffected by the coupling.
\end{enumerate}
\begin{enumerate}[label=\textup{(Z5$'$)},leftmargin=12mm]
  \setcounter{enumi}{4}
  \item \label{zone:Z5prime}
  $S(N)\supseteq\Zseed(N)$.
\end{enumerate}
With the uniform assignment $r_N(y)=B(N)$, the last clause becomes
\[
 S(N)\supseteq\Legal(P_N)\cap B_{8(B(N)-1)}
 \bigl(\Prot(N)\setminus(\Legal(P_N)\cup\Stones(P_N))\bigr).
\]
All four components are finite and verifier-enumerable.  The virgin query can
be inverted from each legal candidate over the finite radius allowed by
$E_N^\D(W)\le B(N)$.
\end{definition}

\begin{plainterms}
Direct obligations are legal cells that the defender could occupy now.  Seed
guards stop a sequence of dismissed moves from carrying legality toward an
obligation that is not yet legal.  Touched and virgin guards stop the same
kind of sequence from completing a defender window.  Current threat-hitting
cells are useful search candidates, but they are not a premise of this zone
theorem.
\end{plainterms}

The seed radius is set by the number of remaining opportunities to advance a
radius-$8$ legality chain.

\begin{lemma}[First protected occupation]
\label{lem:L9prime}
Suppose every direct legal protected cell is searched and every ghost-legal
dismissal $x$ at a node $N$ lies outside radius $8(B(N)-1)$ of
\[
  \Prot(N)\setminus(\Legal(P_N)\cup\Stones(P_N)).
\]
Then no defender placement creates a real-only stone in the current protected
set.  In the ranked form, use the maximum live role rank $r_N(y)$ for each
target $y$ in place of $B(N)$; the same conclusion holds through that role's
deadline.
\end{lemma}

\begin{proof}[Proof sketch]
Assume the first bad occupation is at $y$.  If $y$ is ghost-legal, it lies in
$\Zdir$ and is copied.  Otherwise trace real-only legality backward to the
first ghost-legal dismissed seed $x_0$; if $y$ is the $p$th defender placement
in that chain, then $\dist(x_0,y)\le8(p-1)$ and nesting keeps $y$ protected at
$x_0$.  The local budget gives $p\le B(N_{x_0})$, or the live role gives
$p\le r_{N_{x_0}}(y)$, so $x_0$ was in the required seed band.  See
\cite{hexobot-companions}.
\end{proof}

The distance inequality is attained for an exact rank $r$: a ghost-legal seed
may be followed by $r-1$ successive distance-$8$ defender placements, with
the protected target last.  This does not establish sharpness of the uniform
$B$-only wrapper.

\subsection{The real--ghost coupling}

The coupling follows an exact certificate path even when the real defender
plays an omitted reply.  Canonical set differences record the defender stones
that occur in only one of the two positions.

\begin{definition}[Real--ghost coupling]
\label{def:D12}
At each coupled step, the real position $R$ and ghost position $G$ have the
same ply, mover, and remaining budget, and have identical attacker stones.
Define
\[
 X=\Stones_\D(R)\setminus\Stones_\D(G),
 \qquad
 Y=\Stones_\D(G)\setminus\Stones_\D(R).
\]
Every member of $X$ entered at a dismissal for which the ordinary zone or a
selected \cref{def:D17} envelope certified avoidance of the live obligations.
Every member of $Y$ entered as a searched ghost filler.  The coupling
maintains
\[
  X\cap\Prot(N)=\varnothing
\]
at every visited node.  A history $\widehat X\supseteq X$ retains every cell
ever added to $X$.  For every window $W$ the canonical identity and its useful
upper bound are
\begin{equation}
\cnt_\D(W,R)=\cnt_\D(W,G)+|X\cap W|-|Y\cap W|
 \le\cnt_\D(W,G)+|\widehat X\cap W|.
\tag{MI}\label{eq:MI}
\end{equation}
\end{definition}

\begin{figure}[tbp]
  \centering
  \begin{tikzpicture}[scale=.38]
    \begin{scope}[xshift=-8.1cm]
      \HexBoardRange{0}{0}{3}
      \HexAttacker{-1}{0}\HexAttacker{0}{0}\HexAttacker{1}{-1}
      \HexDefender{-1}{1}\HexDefender{2}{-1}
      \node[defender stone,fill=black!45] at \HexAt{1}{1} {};
      \HexMark{1}{1}{cross}{}
      \node[font=\scriptsize,above] at \HexAt{1}{1} {$X$};
      \node[font=\small] at (0,-6.8) {real position $R$};
    \end{scope}
    \begin{scope}[xshift=8.1cm]
      \HexBoardRange{0}{0}{3}
      \HexAttacker{-1}{0}\HexAttacker{0}{0}\HexAttacker{1}{-1}
      \HexDefender{-1}{1}\HexDefender{2}{-1}\HexDefender{-2}{0}
      \node[font=\scriptsize,above] at \HexAt{-2}{0} {$Y$};
      \node[font=\small] at (0,-6.8) {ghost position $G$};
    \end{scope}
  \end{tikzpicture}
  \caption{A schematic coupled pair.  Attacker stones and the two unmarked
  defender stones agree.  The crossed defender stone belongs only to $R$ and
  hence to $X$; the labelled open stone belongs only to $G$ and hence to $Y$.
  Both boards retain the same mover, budget, and ply.}
  \label{fig:ghost-coupling}
\end{figure}

\begin{plainterms}
The ghost game consumes one exact certificate edge on every defender ply.
When the real defender chooses another cell, the two defender boards may
differ, but their attacker board and turn state remain synchronized.  The
sets $X$ and $Y$ make every occupancy and mask discrepancy explicit.
\end{plainterms}

The completion guard uses the mask inequality together with a separate
causal chain for the first omitted fill of a window.

\begin{lemma}[Split completion safety]
\label{lem:L12}
In the coupling of \cref{def:D12}, under \cref{zone:Z2,zone:Z4} and
\hyperref[zone:Z5prime]{clause \textup{(Z5$'$)}}, no real defender play
completes a fixed window $W$ before the mapped certificate attacker wins or
first enters $W$.
\end{lemma}

\begin{proof}[Proof sketch]
If no cell of $W$ was dismissed, \eqref{eq:MI} gives the ghost at least the
real defender count, contradicting the ban on defender-terminal certificate
edges.  If the first omitted fill enters a touched ghost window, the
$\cnt_\D+E^\D\ge6$ test places that fill in $\Ztouch$.  If the ghost window is
virgin, a ghost-legal first fill lies in $\Zvirgin$ directly; a ghost-illegal
fill traces back through $j$ real-only legality links to a seed at distance at
most $8j$, while completion requires exposure at least $j+6$.  The same
count includes every LOSS remainder, and the mandatory leaf
$\neg\ownwinnow$ check bounds a no-dismissal remainder by five stones.  See
\cite{hexobot-companions}.
\end{proof}

For one fixed window the virgin arithmetic is attained at exposure $7$: a
legal seed at distance $8$ may be followed by the first window fill and the
five remaining fills.  This fixed-window trace does not establish sharpness
of the complete union $\Zvirgin$.
